\subsection{Messverfahren}
Für alle Messungen wird ein Gitterspektrometer verwendet, welches das entsprechende Licht über ein Objektiv mit Lichtleitfaser aufnimmt und es durch die dispersiven Eigenschaften eines Gitters so aufspaltet, dass mithilfe eines LCD-Sensors die einzelnen Wellenlängen detektiert werden können. Das resultierende Spektrum (Intensität aufgetragen über Wellenlänge) wird am Computer durch OceanView als Diagramm angezeigt und kann sowohl als \verb|.pdf| als auch als \verb|.txt| Datei abgespeichert werden. Durch Fokussieren und Ausrichten des Objektives auf die jeweilig zu untersuchende Lichtquelle werden die Messergebnisse erhoben.\\
Neben Sonnenlicht, blauem Himmel vor und hinter Glas, sowie den verschiedenen Lichtquellen, wurde die Natriumlampe genauer untersucht. Dafür wurde das Spektrum im Bereich von \(\qtyrange{400}{540}{\nano\meter}\) einmal aufgenommen, sodass der Peak bei \(\qty{500}{\nano\meter}\) in Sättigung ist (also nicht vollständig auf dem Diagramm zu sehen, da die Skalierung der Intensitätsachse zu klein ist), dadurch können die schwachen Peaks besser anvisiert werden, und einmal nicht mit dem Peak in Sättigung.\\
In Aufgabe 4 werden nun zwei Linien bei \(\qty{570}{\nano\meter}\) und \(\qty{615}{\nano\meter}\) sichtbar eingestellt, die sonst von der D-Doppellinie (Dublett) überstrahlt werden, die D-Linie wird auch aufgezeichnet. \\
Aufgabe 5 verlangt nun noch die Aufzeichnung des Spektrums im Bereich von \(\qtyrange{650}{850}{\nano\meter}\).
\subsection{Messprotoll}
Es liegt kein Messprotokoll vor, die erhobenen Daten wurden in den entsprechenden \verb|.txt| Dateien abgespeichert. 
\begin{figure}[htb]
    \centering
    \includegraphics[width=7cm]{redacted.png}
    \caption{These files have been redacted}
\end{figure}
